% Auto-generated — in master document use: \graphicspath{{paper_outputs/<run>/figures/}}

\begin{figure}[t]
  \centering
  \includegraphics[width=0.9\linewidth]{fig\_energy.pdf}
  \caption{Energy relaxation during the SRG simulation.}
  \label{fig:energy}
\end{figure}

\begin{figure}[t]
  \centering
  \includegraphics[width=0.9\linewidth]{fig\_reward.pdf}
  \caption{Reward trajectory including spectral and braid shaping.}
  \label{fig:reward}
\end{figure}

\begin{figure}[t]
  \centering
  \includegraphics[width=0.9\linewidth]{fig\_absorbance.pdf}
  \caption{Absorbance proxy versus step.}
  \label{fig:absorbance}
\end{figure}

\begin{figure}[t]
  \centering
  \includegraphics[width=0.9\linewidth]{fig\_electrolyte.pdf}
  \caption{Ionic strength and Debye length.}
  \label{fig:electrolyte}
\end{figure}

\begin{figure}[t]
  \centering
  \includegraphics[width=0.9\linewidth]{fig\_activity.pdf}
  \caption{Activity coefficients from the selected electrolyte model.}
  \label{fig:activity}
\end{figure}

\begin{figure}[t]
  \centering
  \includegraphics[width=0.9\linewidth]{fig\_conductivity.pdf}
  \caption{Specific conductivity proxy (Onsager/Kohlrausch style).}
  \label{fig:conductivity}
\end{figure}

\begin{figure}[t]
  \centering
  \includegraphics[width=0.9\linewidth]{fig\_spectral.pdf}
  \caption{Spectral gaps and NCG smoothness.}
  \label{fig:spectral}
\end{figure}

\begin{figure}[t]
  \centering
  \includegraphics[width=0.9\linewidth]{fig\_braid.pdf}
  \caption{Braid complexity metrics.}
  \label{fig:braid}
\end{figure}

\begin{figure}[t]
  \centering
  \includegraphics[width=0.9\linewidth]{fig\_performance.pdf}
  \caption{Timing and throughput.}
  \label{fig:performance}
\end{figure}

\begin{figure}[t]
  \centering
  \includegraphics[width=0.9\linewidth]{fig\_dashboard\_summary.pdf}
  \caption{Four-panel overview of key observables.}
  \label{fig:dashboard-summary}
\end{figure}
